% !TEX root = creche.tex
\documentclass[creche.tex]{subfiles}
\chapter{Topological dynamics}\label{dynamics}


\section{}


Let $G$ be an infinite group that acts on $\U^{|{\mr x}|}$. We say that $\grD$ \emph{is quasi-invariant\/} if for every finitely many $f_1,\dots,f_n\in G$ the sets  $f_i[\grD]$ have non-empty intersection. We say that a \emph{formula\/} or a \emph{type is drifting\/} or quasi-invariant if the set it defines is. 


We say that $\grD$ \emph{is drifting\/} if for every finitely many $f_1,\dots,f_n\in G$ there is a $g\in G$ such that $g[\grD]$ is disjoint from all the $f_i[\grD]$. We say that $\grD$ \emph{is quasi-invariant\/} if for every finitely many $f_1,\dots,f_n\in G$ the sets  $f_i[\grD]$ have non-empty intersection. We say that a \emph{formula\/} or a \emph{type is drifting\/} or quasi-invariant if the set it defines is.



The following notions apply generally to any group $G$ acting on some set ${\gr\X}$ and to any set $\grD\subseteq {\gr\X}$. Below we always have $G=\Autf(\U/A)$ and ${\gr\X}=\U^{|{\gr z}|}$.  We say that $\grD$ \emph{is drifting\/} if for every finitely many $f_1,\dots,f_n\in G$ there is a $g\in G$ such that $g[\grD]$ is disjoint from all the $f_i[\grD]$. 



which we assume to be a structure in a language $L$ that extends the language of multiplicative groups with a unary relation for every subset of $G$. If $X\subseteq G$ we write $\xi_X(x)$ for the corresponding predicate. We assume that $\U$ is a large saturated elementary extension of $G$. For $\phi(x)\in L(\U)$ we write $\phi(G)$ for $G\cap\phi(\U)$.
